% Figure: compressibility factor Z against reduced pressure p_r along two
% reduced isotherms (T_r = 1.0 and T_r = 1.5), illustrating the principle
% of corresponding states. Three real gases (nitrogen, methane, carbon
% dioxide) are plotted as points reduced by their own critical constants;
% they collapse onto a near-common curve, the empirical content of the
% corresponding-states principle. The horizontal line marks the van der
% Waals universal critical compressibility Z_c = 3/8 = 0.375, above the
% measured cluster near 0.27 to 0.29, the gap the acentric factor closes.
\begin{figure}[htbp]
\centering
\begin{tikzpicture}
\begin{axis}[
  width=12cm, height=8cm,
  xlabel={reduced pressure $p_r = p/p_c$},
  ylabel={compressibility factor $Z = pv/(R_s T)$},
  xmin=0, xmax=4, ymin=0.2, ymax=1.1,
  axis lines=left,
  tick label style={font=\scriptsize},
  label style={font=\small},
  legend pos=south east,
  legend style={font=\scriptsize},
]
% generalised corresponding-states curve at T_r = 1.0 (dips deepest)
\addplot[engblue, very thick, smooth] coordinates {
  (0,1.00) (0.5,0.86) (1.0,0.71) (1.5,0.55) (2.0,0.42)
  (2.5,0.40) (3.0,0.44) (3.5,0.50) (4.0,0.57)
};
\addlegendentry{$T_r = 1.0$ (corresponding states)}
% generalised curve at T_r = 1.5 (shallow dip, returns above 1)
\addplot[engred, very thick, smooth] coordinates {
  (0,1.00) (0.5,0.95) (1.0,0.90) (1.5,0.87) (2.0,0.86)
  (2.5,0.87) (3.0,0.90) (3.5,0.94) (4.0,0.99)
};
\addlegendentry{$T_r = 1.5$ (corresponding states)}
% three real gases reduced by their own critical constants, T_r = 1.0
% landing near the blue curve to show the collapse
\addplot[only marks, mark=*, engblack, mark size=2.0pt] coordinates {
  (0.5,0.85) (1.0,0.72) (2.0,0.43)
};
\addlegendentry{N$_2$, CH$_4$, CO$_2$ reduced ($T_r=1.0$)}
% van der Waals universal Z_c = 3/8
\addplot[engorange, very thick, dashed, domain=0:4] {0.375};
\node[engorange, font=\scriptsize, anchor=south west] at (axis cs:0.1,0.385)
  {van der Waals $Z_c = 3/8 = 0.375$};
% measured Z_c band near 0.27 to 0.29
\addplot[enggray, thick, dotted, domain=0:4] {0.28};
\node[enggray, font=\scriptsize, anchor=north west] at (axis cs:0.1,0.27)
  {measured $Z_c \approx 0.27$--$0.29$};
\end{axis}
\end{tikzpicture}
\caption[Corresponding-states compressibility chart]{Compressibility
factor against reduced pressure along two reduced isotherms. When each
gas is plotted in its own reduced coordinates $(p_r, T_r)$, the data for
nitrogen, methane, and carbon dioxide collapse onto a near-common curve
(black points falling on the $T_r = 1.0$ line), the empirical content of
the principle of corresponding states: gases at the same fraction of
their critical temperature and pressure deviate from ideality by nearly
the same factor. The curve at $T_r = 1.0$ plunges to $Z \approx 0.4$ near
$p_r = 2$ to $2.5$, while the $T_r = 1.5$ curve barely dips, the working
basis of the rule that ideality is safe at $T_r > 2$. The dashed line is
the van der Waals universal critical compressibility $Z_c = 3/8$, which
overshoots the measured cluster near $0.27$ to $0.29$ (dotted line); the
gap is what the acentric factor and the cubic equations of state are
fitted to close.}
\label{fig:vol04ch02:zchart-corresponding-states}
\end{figure}
